% AA-12 | Algebra abstrakcyjna | Iteracja VI
\chapter{Pierwiastniki ideałów i lokalizacja pierścieni}
\chaptermark{Pierwiastniki ideałów i lokalizacja}
\label{aa:ch:lokalizacja}
\setlength{\parskip}{2pt}

\section{Dlaczego zbiór rozwiązań nie wyznacza ideału?}
W Rozdziale~\ref{aa:ch:wielomiany} porównaliśmy ideały
$(x)$ oraz $(x^2)$ w $\mathbb R[x]$. Oba opisują ten sam
zbiór zer, ale tylko drugi iloraz ma niezerowy element
nilpotentny. Pierwiastnik ideału pozwala oddzielić część
równania widoczną przez zwykłe zerowanie od dodatkowych
informacji algebraicznych. Lokalizacja odpowiada na inne
pytanie: jakie ułamki wolno dołączyć do pierścienia, jeśli
chcemy badać zachowanie algebraiczne w wybranym miejscu?

W całym rozdziale, jeśli nie podano inaczej, $R$ jest
\emph{przemiennym pierścieniem z jedynką}, a ideały
są dwustronne. Dopuszczamy pierścień zerowy tylko tam,
gdzie jest to jawnie dozwolone.

\section{Pierwiastnik ideału i nilpotentność}
\begin{definicja}[Pierwiastnik ideału]
\label{aa:def:ideal-radical}
Niech $R$ będzie przemiennym pierścieniem z jedynką,
a $I\subseteq R$ jego ideałem. \emph{Pierwiastnikiem}
ideału $I$ nazywamy zbiór
\[
 \sqrt I=\{a\in R:\ \exists n\in\mathbb N,\ n\ge1,\ a^n\in I\}.
\]
Ideał $I$ nazywamy \emph{radykalnym}, jeśli $\sqrt I=I$.
Pierwiastnik ideału zerowego $\sqrt{(0)}$ nazywamy
\emph{nilradykałem} pierścienia $R$.
\end{definicja}

\begin{twierdzenie}[Pierwiastnik ideału jest ideałem]
\label{aa:tw:radical-ideal}
Niech $R$ będzie przemiennym pierścieniem z jedynką,
a $I\subseteq R$ jego ideałem. Wówczas zbiór $\sqrt I$
jest ideałem zawierającym $I$.
Ponadto $\sqrt{\sqrt I}=\sqrt I$, a dla każdego
ideału $J\subseteq R$ z $I\subseteq J$ zachodzi
$\sqrt I\subseteq\sqrt J$.
\end{twierdzenie}
\begin{proof}
Jeśli $a\in I$, to $a^1\in I$, więc $I\subseteq\sqrt I$.
Dla $a\in\sqrt I$ i $r\in R$ istnieje $m\ge1$ z
$a^m\in I$. Wtedy $(ra)^m=r^ma^m\in I$,
więc $ra\in\sqrt I$. Również $-a\in\sqrt I$, gdyż
$(-a)^m=(-1)^ma^m\in I$.

Pozostaje dowieść zamkniętości na dodawanie.
Niech $a^m\in I$ i $b^n\in I$ dla pewnych $m,n\ge1$.
Rozwińmy $(a+b)^{m+n-1}$ za pomocą dwumianu Newtona.
W każdym składniku $\binom{m+n-1}{k}a^k b^{m+n-1-k}$
albo $k\ge m$, albo $m+n-1-k\ge n$: w przeciwnym
razie mielibyśmy $k\le m-1$ i
$m+n-1-k\le n-1$, co po dodaniu daje sprzeczność.
Pierwszy rodzaj składników należy do $I$, bo jest
podzielny przez $a^m$, a drugi należy do $I$, bo jest
podzielny przez $b^n$. Cała suma należy do $I$.
Wobec tego $a+b\in\sqrt I$ i $\sqrt I$ jest ideałem.

Jeśli $a\in\sqrt{\sqrt I}$, to $a^m\in\sqrt I$
dla pewnego $m\ge1$, a następnie $(a^m)^n\in I$
dla pewnego $n\ge1$. Zatem $a^{mn}\in I$,
czyli $a\in\sqrt I$. Odwrotne zawieranie wynika
z $\sqrt I\subseteq\sqrt{\sqrt I}$. Monotoniczność
jest bezpośrednia: $a^n\in I\subseteq J$
pociąga $a\in\sqrt J$.
\end{proof}

\begin{twierdzenie}[Pierwiastnik iloczynu i przecięcia ideałów]
\label{aa:tw:radical-product}
Niech $R$ będzie przemiennym pierścieniem z jedynką,
a $I,J\subseteq R$ jego ideałami. Wtedy zachodzi
$\sqrt{IJ}=\sqrt{I\cap J}=\sqrt I\cap\sqrt J$.
\end{twierdzenie}
\begin{proof}
Ponieważ $IJ\subseteq I\cap J$, z monotoniczności
pierwiastnika otrzymujemy
$\sqrt{IJ}\subseteq\sqrt{I\cap J}\subseteq\sqrt I\cap\sqrt J$.
Jeśli $a\in\sqrt I\cap\sqrt J$, to istnieją
$m,n\ge1$ takie, że $a^m\in I$ oraz $a^n\in J$.
Wówczas $a^{m+n}=a^m a^n\in IJ$. Oznacza to
$a\in\sqrt{IJ}$, co daje brakujące zawieranie
i wszystkie równości.
\end{proof}

\begin{stwierdzenie}[Nilradykał i pierścienie zredukowane]
\label{aa:st:nilradical-reduced}
Niech $R$ będzie przemiennym pierścieniem z jedynką,
$I\subseteq R$ jego ideałem, a $a\in R$. Wówczas element
$a+I\in R/I$ jest nilpotentny wtedy i tylko wtedy,
gdy $a\in\sqrt I$. W szczególności
$R/I$ nie ma niezerowych elementów nilpotentnych wtedy
i tylko wtedy, gdy $I$ jest radykalny.
Pierścień bez niezerowych elementów nilpotentnych
nazywamy \emph{pierścieniem zredukowanym}.
\end{stwierdzenie}
\begin{proof}
Z definicji działań w ilorazie
$(a+I)^n=a^n+I$. Jest to klasa zerowa dokładnie wtedy,
gdy $a^n\in I$. Istnienie takiego $n\ge1$ jest
równoważne $a\in\sqrt I$. Jeśli wszystkie elementy
nilpotentne ilorazu są zerowe, to każdy $a\in\sqrt I$
należy do $I$, zatem $\sqrt I=I$. Odwrotnie,
gdy $I=\sqrt I$, dowolny element nilpotentny
$a+I$ ma reprezentanta $a\in I$, więc jest zerowy.
\end{proof}

\begin{przyklad}[Trzy ideały opisujące podobne równania]
\label{aa:prz:radical-examples}
Niech $K$ będzie ciałem. W $K[x]$ mamy
$\sqrt{(x^m)}=(x)$ dla każdego $m\ge1$.
Istotnie, $x^m\in(x^m)$, więc $x\in\sqrt{(x^m)}$
i $(x)\subseteq\sqrt{(x^m)}$.
Z drugiej strony $(x^m)\subseteq(x)$, a $(x)$ jest
ideałem pierwszym: $K[x]/(x)\cong K$ jest ciałem.
Ideał pierwszy jest radykalny, ponieważ
z $a^n\in(x)$ jego pierwszość daje kolejno $a\in(x)$.
Stąd $\sqrt{(x^m)}\subseteq(x)$.

W $K[x,y]$ rozważmy $I=(x^2,xy)$. Zachodzi
$I\subseteq(x)$ i $x^2\in I$. Prowadząc dokładnie ten
sam argument, otrzymujemy $\sqrt I=(x)$.
Jednak $I\ne(x)$: gdyby $x=x^2F+xyG$, to
po wyłączeniu $x$ i skróceniu w dziedzinie całkowitości
mielibyśmy $1=xF+yG$, co jest niemożliwe po
podstawieniu $x=y=0$. W szczególności ideał $I$
nie jest radykalny.
\end{przyklad}

\begin{przyklad}[Nilradykał pierścienia $\mathbb Z/12\mathbb Z$]
\label{aa:prz:nilradical-z12}
W pierścieniu $R=\mathbb Z/12\mathbb Z$ klasa $[6]$
jest niezerowym elementem nilpotentnym, gdyż
$[6]^2=[36]=[0]$. Dokładnie elementy ideału $([6])$
są nilpotentne. Istotnie, jeśli $[a]^n=[0]$, to
$12\mid a^n$, a więc $2\mid a$ i $3\mid a$.
Stąd $6\mid a$ i $[a]\in([6])$. Odwrotnie,
$([6])=\{[0],[6]\}$ i obie klasy są nilpotentne.
Ponadto $([2])$ i $([3])$ są ideałami pierwszymi,
bo ich ilorazy są odpowiednio izomorficzne z
$\mathbb F_2$ i $\mathbb F_3$; ich przecięcie wynosi
$([6])$. Są to wszystkie ideały pierwsze $R$:
jeśli $P$ jest ideałem pierwszym, to
$[2]^2[3]=[12]=[0]\in P$, więc $[2]\in P$
lub $[3]\in P$. W pierwszym przypadku
$([2])\subseteq P$, a w drugim $([3])\subseteq P$;
ponieważ oba te ideały są maksymalne i $P\ne R$,
mamy odpowiednio $P=([2])$ lub $P=([3])$.
Otrzymujemy zatem konkretną tożsamość
$\sqrt{(0)}=\bigcap_{P\ \mathrm{pierwszy}}P=([6])$.
\end{przyklad}

\section{Pierwiastnik jako przecięcie ideałów pierwszych}
W następnym twierdzeniu użyjemy lematu Kuratowskiego--Zorna:
jeżeli w niepustym zbiorze częściowo uporządkowanym każdy
łańcuch ma ograniczenie górne, to istnieje w nim element
maksymalny. Potrzebny nam argument jest konstrukcją
maksymalnego ideału unikającego wskazanych potęg elementu.

\begin{lemat}[Ideał pierwszy rozłączny ze zbiorem multiplikatywnym]
\label{aa:lem:prime-avoid-S}
Niech $R$ będzie przemiennym pierścieniem z jedynką,
$I\subseteq R$ ideałem, a $S\subseteq R$ niepustym zbiorem
zawierającym $1_R$ i zamkniętym na mnożenie.
Jeżeli $I\cap S=\varnothing$, to istnieje ideał
pierwszy $P$ taki, że $I\subseteq P$ oraz
$P\cap S=\varnothing$.
\end{lemat}
\begin{proof}
Rozważmy rodzinę $\mathcal F$ wszystkich ideałów
$J$ z $I\subseteq J$ i $J\cap S=\varnothing$,
uporządkowaną przez zawieranie. Rodzina jest niepusta,
bo $I\in\mathcal F$. Pusty łańcuch ma ograniczenie
górne $I$. Jeśli natomiast $\mathcal C\subseteq\mathcal F$
jest niepustym łańcuchem, to suma mnogościowa
$\bigcup_{J\in\mathcal C}J$ jest ideałem: dla dwóch jej elementów wybieramy ideał
z łańcucha zawierający oba, co jest możliwe dzięki
liniowemu uporządkowaniu łańcucha; w nim sprawdzamy
zamkniętość na odejmowanie i mnożenie przez pierścień.
Suma zawiera $I$ i nie przecina $S$, bo żaden z jej
składników nie przecina $S$. Jest zatem ograniczeniem
górnym dla łańcucha. Lemat Kuratowskiego--Zorna daje
maksymalny element $P$ rodziny $\mathcal F$.

Pokażemy, że $P$ jest pierwszy. Jest właściwy, ponieważ
$1_R\in S$, a $P\cap S=\varnothing$.
Załóżmy przeciwnie, że $ab\in P$, lecz $a\notin P$
i $b\notin P$. Ideały $P+(a)$ i $P+(b)$ ściśle
zawierają $P$. Na mocy maksymalności $P$ żaden
z nich nie może należeć do $\mathcal F$.
Istnieją więc $s\in(P+(a))\cap S$ oraz
$t\in(P+(b))\cap S$. Zapiszmy
$s=p+ra$, $t=q+ub$ dla $p,q\in P$, $r,u\in R$.
Wtedy $st=pq+pub+qra+ruab\in P$, ponieważ
pierwsze trzy składniki zawierają czynnik z $P$,
a ostatni należy do $P$ z założenia $ab\in P$.
Jednak $st\in S$ dzięki zamkniętości $S$ na mnożenie.
Otrzymaliśmy sprzeczność z $P\cap S=\varnothing$.
Zatem $P$ jest ideałem pierwszym.
\end{proof}

\begin{twierdzenie}[Pierwiastnik ideału jako przecięcie ideałów pierwszych]
\label{aa:tw:radical-intersection-primes}
Niech $R$ będzie przemiennym pierścieniem z jedynką
i niech $I$ będzie jego ideałem. Wtedy
\[
 \sqrt I=\bigcap_{\substack{P\text{ -- ideał pierwszy}\\ I\subseteq P}}P.
\]
Jeżeli nie ma żadnego ideału pierwszego zawierającego
$I$, przecięcie po pustej rodzinie oznacza cały $R$.
\end{twierdzenie}
\begin{proof}
Jeżeli $I=R$, obie strony są równe $R$ na mocy
przyjętej konwencji. Załóżmy więc $I\ne R$.
Niech $a\in\sqrt I$; dla pewnego $n\ge1$ mamy
$a^n\in I$. Jeśli $P$ jest ideałem pierwszym zawierającym
$I$, to $a^n\in P$. Wielokrotne zastosowanie własności
ideału pierwszego daje $a\in P$. Zatem $a$ należy
do wszystkich takich $P$.

Dla odwrotnego zawierania weźmy $a\notin\sqrt I$.
Wtedy żadna dodatnia potęga $a$ nie należy do $I$;
ponieważ $I$ jest właściwy, również $1_R\notin I$.
Zbiór $S=\{1_R,a,a^2,\ldots\}$ jest multiplikatywny
i nie przecina $I$. Poprzedni lemat daje ideał
pierwszy $P\supseteq I$ rozłączny z $S$.
W szczególności $a\notin P$. Zatem $a$ nie należy
do przecięcia wszystkich ideałów pierwszych zawierających
$I$. Otrzymujemy wymaganą równość.
\end{proof}

\begin{uwaga}[Interpretacja geometryczna]
\label{aa:uw:radical-geometry}
Jeżeli $f\in K[x_1,\ldots,x_n]$ i $f^m$ zanika
w punkcie $z\in K^n$, to także $f$ zanika w tym punkcie,
gdyż w ciele nie ma niezerowych elementów nilpotentnych.
Dlatego ideały $I$ i $\sqrt I$ mają ten sam zbiór
wspólnych zer. Nie twierdzimy tu jednak, że każdy
wielomian znikający na takim zbiorze należy do $\sqrt I$:
taka równość wymaga dodatkowych założeń i prowadzi
do twierdzenia Hilberta o zerach.
\end{uwaga}

\begin{przyklad}[Dlaczego ciało algebraicznie domknięte ma znaczenie]
\label{aa:prz:real-zero-set-not-radical}
Dla $K=\mathbb R$ rozważmy $I=(x^2+1)\subset\mathbb R[x]$.
Zbiór rzeczywistych zer jest pusty: dla każdej
liczby rzeczywistej $a$ zachodzi $a^2+1>0$.
Z tego powodu każdy wielomian znika na wszystkich
punktach tego zbioru (warunek jest spełniony
pusto), więc ideał wielomianów znikających na zbiorze
zer $I$ jest całym $\mathbb R[x]$.
Tymczasem $I$ jest ideałem maksymalnym,
ponieważ $\mathbb R[x]/(x^2+1)$ jest ciałem,
izomorficznym z $\mathbb C$. Ideał maksymalny
jest pierwszy, a każdy ideał pierwszy jest radykalny;
zatem $\sqrt I=I\ne\mathbb R[x]$.
Sama znajomość rzeczywistych zer nie pozwala więc
odzyskać pierwiastnika dowolnego ideału.
\end{przyklad}

\section{Lokalizacja: konstrukcja przez pary i usuwanie mianowników}
\begin{definicja}[Zbiór multiplikatywny]
\label{aa:def:multiplicative-set}
Niech $R$ będzie przemiennym pierścieniem z jedynką.
Podzbiór $S\subseteq R$ jest \emph{multiplikatywny},
jeśli $1_R\in S$ oraz $st\in S$ dla wszystkich $s,t\in S$.
W konstrukcji niezerowej lokalizacji zakładamy dodatkowo
$0_R\notin S$. Gdyby $0_R\in S$, lokalizacja byłaby
pierścieniem zerowym.
\end{definicja}

Nie wystarczy przyjąć zwykłej relacji $as'=a's$
na parach $(a,s)$. W pierścieniach z dzielnikami zera
tak zdefiniowana relacja nie musi być przechodnia.
Należy dopuścić dodatkowy czynnik należący do $S$.

\begin{przyklad}[Zawodność naiwnej relacji na ułamkach]
\label{aa:prz:naive-local-relation}
W $R=\mathbb Z/6\mathbb Z$ weźmy
$S=\{[1],[2],[4]\}$; jest to zbiór multiplikatywny
niezawierający zera. W relacji zadanej wyłącznie
równością iloczynów krzyżowych mamy
$( [0],[1])\sim([0],[2])$, ponieważ $0=0$,
oraz $([0],[2])\sim([3],[1])$, ponieważ
$[0]\cdot[1]=[3]\cdot[2]=[0]$.
Jednak $([0],[1])\not\sim([3],[1])$, gdyż
$[0]\cdot[1]\ne[3]\cdot[1]$.
Relacja nie jest więc przechodnia. W poprawnej
konstrukcji mnożnik $[2]\in S$ zeruje różnicę
$[0]\cdot[1]-[3]\cdot[1]=[3]$.
\end{przyklad}

\begin{definicja}[Lokalizacja pierścienia]
\label{aa:def:localization}
Niech $R$ będzie przemiennym pierścieniem z jedynką,
a $S\subseteq R$ zbiorem multiplikatywnym niezawierającym zera.
Na $R\times S$ definiujemy relację
\[
 (a,s)\sim(b,t)
 \quad\Longleftrightarrow\quad
 \exists u\in S:\ u(at-bs)=0_R.
\]
Klasę pary $(a,s)$ oznaczamy $a/s$. Zbiór klas
nazywamy \emph{lokalizacją} $R$ względem $S$ i oznaczamy
$S^{-1}R$. Na klasach określamy działania
\[
 \frac as+\frac bt=\frac{at+bs}{st},\qquad
 \frac as\cdot\frac bt=\frac{ab}{st}.
\]
\end{definicja}

\begin{twierdzenie}[Poprawność konstrukcji lokalizacji]
\label{aa:tw:localization-well-defined}
Niech $R$ będzie przemiennym pierścieniem z jedynką,
a $S\subseteq R$ zbiorem multiplikatywnym niezawierającym zera.
Dla $(a,s),(b,t)\in R\times S$ określmy
$(a,s)\sim(b,t)$ wtedy i tylko wtedy, gdy istnieje $u\in S$
takie, że $u(at-bs)=0_R$. Wówczas $\sim$ jest relacją
równoważności na $R\times S$, a wzory
\[
\frac as+\frac bt=\frac{at+bs}{st},\qquad
\frac as\cdot\frac bt=\frac{ab}{st}
\]
są dobrze określone na klasach i czynią $S^{-1}R$
przemiennym pierścieniem z jedynką. Odwzorowanie
$\iota:R\to S^{-1}R$, $a\mapsto a/1_R$, jest
homomorfizmem pierścieni; każdy $s/1_R$ dla $s\in S$
jest jednostką lokalizacji o odwrotności $1_R/s$.
\end{twierdzenie}
\begin{proof}
\emph{Relacja równoważności.} Zwrotność wynika z
$1_R(as-sa)=0_R$, a symetria z zamiany znaku
różnicy. Dla przechodniości załóżmy
$(a,s)\sim(b,t)$ oraz $(b,t)\sim(c,v)$.
Istnieją $u,w\in S$ takie, że
$u(at-bs)=0_R$ i $w(bv-ct)=0_R$.
Pomnóżmy pierwszą równość przez $wv$,
a drugą przez $us$. Po dodaniu otrzymujemy
\[
 uwt(av-cs)=0_R.
\]
Ponieważ $uwt\in S$, zachodzi $(a,s)\sim(c,v)$.

\emph{Niezależność działań od reprezentantów.}
Załóżmy $(a,s)\sim(a',s')$ i $(b,t)\sim(b',t')$.
Wtedy pewne $u,w\in S$ zerują odpowiednio
$as'-a's$ oraz $bt'-b't$. Dla dodawania różnica
liczników po sprowadzeniu wyników do wspólnego mianownika
jest równa
\[
 (at+bs)s't'-(a't'+b's')st
 =tt'(as'-a's)+ss'(bt'-b't).
\]
Po pomnożeniu przez $uw\in S$ oba składniki
po prawej znikają, więc sumy wyznaczają tę samą klasę.
Podobnie dla mnożenia różnica wynosi
\[
 ab s't'-a'b' st
 =bt'(as'-a's)+a's(bt'-b't),
\]
która również znika po pomnożeniu przez $uw$.
Mianowniki wynikowych ułamków należą do $S$.

\emph{Aksjomaty pierścienia.} Łączność i przemienność
obu działań oraz rozdzielność wynikają z odpowiednich
tożsamości w $R$ po sprowadzeniu ułamków do wspólnych
mianowników. Dla przykładu
$((a/s)+(b/t))+(c/v)=(atv+bsv+cst)/(stv)$,
co nie zależy od nawiasowania. Elementem zerowym
jest $0_R/1_R$, elementem przeciwnym do $a/s$
jest $(-a)/s$, a jedynką $1_R/1_R$.

Wzory na działania pokazują, że $\iota$ zachowuje
dodawanie, mnożenie i jedynkę. Dla $s\in S$ mamy
$(s/1_R)(1_R/s)=s/s=1_R/1_R$,
ponieważ $s\cdot1_R-1_R\cdot s=0_R$.
Stąd $s/1_R$ jest odwracalny.
\end{proof}

\begin{stwierdzenie}[Kiedy element staje się zerem?]
\label{aa:st:local-zero-kernel}
Niech $R$ będzie przemiennym pierścieniem z jedynką,
a $S\subseteq R$ zbiorem multiplikatywnym niezawierającym
zera. Niech $\iota:R\to S^{-1}R$ będzie homomorfizmem
kanonicznym, $\iota(a)=a/1_R$. Dla $a\in R$ zachodzi
$\iota(a)=0$ w $S^{-1}R$ wtedy i tylko wtedy,
gdy istnieje $s\in S$ z $sa=0_R$.
W szczególności $\iota$ jest różnowartościowe, jeśli
żaden element $S$ nie jest dzielnikiem zera.
Gdy $R$ jest dziedziną całkowitości, warunek ten zachodzi
automatycznie dla $0_R\notin S$.
\end{stwierdzenie}
\begin{proof}
Równość $a/1_R=0_R/1_R$ oznacza z definicji relacji
istnienie $s\in S$ z
$s(a\cdot1_R-0_R\cdot1_R)=sa=0_R$.
Jeśli żaden element $s\in S$ nie jest dzielnikiem zera,
równość ta wymusza $a=0_R$,
więc jądro $\iota$ jest zerowe.
\end{proof}

\begin{przyklad}[Lokalizacja liczb całkowitych]
\label{aa:prz:local-z}
Dla $R=\mathbb Z$ i $S=\{2^n:n\ge0\}$
lokalizacja jest pierścieniem
$\mathbb Z[1/2]=\{a/2^n:a\in\mathbb Z,\ n\ge0\}\subset\mathbb Q$.
Dołączyliśmy odwrotność liczby $2$, ale nie
odwrotność liczby $3$. Gdyby $1/3=a/2^n$,
to $3a=2^n$, co wymagałoby $3\mid 2^n$,
a to jest niemożliwe.

Przy $S=\{2^a3^b:a,b\in\mathbb N_0\}$ otrzymujemy
$\mathbb Z[1/6]$, a jego jednostkami są dokładnie
liczby $\pm2^m3^n$, gdzie $m,n\in\mathbb Z$.
Rzeczywiście, te liczby mają odwrotności tej samej postaci.
Jeżeli ułamek $u/(2^a3^b)$ jest jednostką, to po
pomnożeniu przez odwrotność również należącą do
$\mathbb Z[1/6]$ jego licznik $u$ musi dzielić pewną
liczbę postaci $2^c3^d$ w $\mathbb Z$. Jednoznaczny
rozkład na czynniki w $\mathbb Z$ wymusza
$u=\pm2^r3^s$ dla pewnych $r,s\ge0$.
\end{przyklad}

\section{Własność uniwersalna lokalizacji i zachowanie ideałów}
\begin{twierdzenie}[Własność uniwersalna lokalizacji]
\label{aa:tw:localization-universal}
Niech $R,T$ będą przemiennymi pierścieniami z jedynką,
$S\subseteq R$ zbiorem multiplikatywnym niezawierającym
zera i $f:R\to T$ homomorfizmem takim, że
$f(s)\in T^\times$ dla każdego $s\in S$.
Niech $\iota:R\to S^{-1}R$ oznacza homomorfizm
kanoniczny $\iota(a)=a/1_R$. Istnieje dokładnie jeden
homomorfizm pierścieni $\bar f:S^{-1}R\to T$
spełniający $\bar f\circ\iota=f$.
Jest zadany wzorem $\bar f(a/s)=f(a)f(s)^{-1}$.
\end{twierdzenie}
\begin{proof}
Najpierw sprawdzamy niezależność od reprezentantów.
Jeśli $a/s=b/t$, to pewne $u\in S$ spełnia
$u(at-bs)=0_R$. Po zastosowaniu $f$ dostajemy
$f(u)(f(a)f(t)-f(b)f(s))=0_T$.
Ponieważ $f(u)$ jest jednostką, możemy pomnożyć
przez $f(u)^{-1}$, a następnie przez
$f(s)^{-1}f(t)^{-1}$, otrzymując
$f(a)f(s)^{-1}=f(b)f(t)^{-1}$.
Zatem wzór ma jednoznaczny sens na klasach.
Zachowanie dodawania i mnożenia wynika ze zwykłych
wzorów na działania na ułamkach; jedynka przechodzi
na $f(1_R)f(1_R)^{-1}=1_T$.

Dla jednoznaczności zauważmy, że w lokalizacji
$a/s=(a/1_R)(s/1_R)^{-1}$.
Każdy homomorfizm przedłużający $f$ musi więc
wysyłać $a/s$ na $f(a)f(s)^{-1}$. Otrzymana mapa
jest jedynym możliwym przedłużeniem.
\end{proof}

\begin{definicja}[Rozszerzenie ideału do lokalizacji]
\label{aa:def:ideal-extension-local}
Niech $R$ będzie przemiennym pierścieniem z jedynką,
$S\subseteq R$ zbiorem multiplikatywnym niezawierającym
zera, a $I\subseteq R$ ideałem. Przez $S^{-1}I$ oznaczamy
zbiór klas $a/s\in S^{-1}R$,
gdzie $a\in I$, $s\in S$. Jest to ideał pierścienia
$S^{-1}R$ i nazywamy go \emph{rozszerzeniem ideału} $I$.
\end{definicja}
\begin{lemat}[Rozszerzenie ideału jest ideałem]
\label{aa:lem:local-ideal-is-ideal}
Niech $R$ będzie przemiennym pierścieniem z jedynką,
$S\subseteq R$ zbiorem multiplikatywnym niezawierającym
zera, a $I\subseteq R$ ideałem. Wtedy zbiór $S^{-1}I$
jest ideałem pierścienia $S^{-1}R$.
\end{lemat}
\begin{proof}
Element $0/1$ należy do $S^{-1}I$, bo $0\in I$.
Jeśli $a/s,b/t\in S^{-1}I$, gdzie $a,b\in I$, to
$(a/s)-(b/t)=(at-bs)/(st)$, a $at-bs\in I$,
ponieważ $I$ jest ideałem. Zatem zbiór jest podgrupą
addytywną. Dla dowolnego $c/u\in S^{-1}R$ mamy
$(c/u)(a/s)=ca/(us)$, a $ca\in I$.
To dowodzi zamkniętości na mnożenie przez elementy lokalizacji.
\end{proof}

\begin{lemat}[Przynależność do rozszerzenia ideału]
\label{aa:lem:local-ideal-membership}
Niech $R$ będzie przemiennym pierścieniem z jedynką,
$S\subseteq R$ zbiorem multiplikatywnym niezawierającym
zera, a $I\subseteq R$ ideałem. Dla $a\in R$ i $s\in S$ zachodzi
$a/s\in S^{-1}I$ wtedy i tylko wtedy, gdy
istnieje $u\in S$ takie, że $ua\in I$.
Ponadto $S^{-1}I$ jest ideałem właściwym wtedy
i tylko wtedy, gdy $I\cap S=\varnothing$.
\end{lemat}
\begin{proof}
Jeżeli $a/s=b/t$ dla $b\in I$, $t\in S$,
to istnieje $v\in S$ z $v(at-bs)=0$.
Wobec tego $(vt)a=(vs)b\in I$, a $vt\in S$.
Odwrotnie, jeśli $ua\in I$, to
$a/s=(ua)/(us)$, ponieważ iloczyny krzyżowe są
równe, zatem ułamek należy do $S^{-1}I$.

Ideał $S^{-1}I$ jest całym pierścieniem dokładnie wtedy,
gdy zawiera jedynkę $1_R/1_R$. Z pierwszej części
jest to równoważne istnieniu $u\in S$ z
$u\cdot1_R=u\in I$, czyli $I\cap S\ne\varnothing$.
\end{proof}

\begin{twierdzenie}[Kontrakcja rozszerzonego ideału i nasycenie]
\label{aa:tw:localization-saturation}
Niech $R$ będzie przemiennym pierścieniem z jedynką,
$S\subseteq R$ zbiorem multiplikatywnym niezawierającym
zera, a $I\subseteq R$ ideałem. Niech
$\iota:R\to S^{-1}R$, $\iota(a)=a/1_R$, będzie
homomorfizmem kanonicznym. Przeciwobraz ideału
$S^{-1}I$ ma postać
\[
 \iota^{-1}(S^{-1}I)=I^S
 :=\{a\in R:\ \exists s\in S\ \ sa\in I\}.
\]
W szczególności $I^S$ jest ideałem zawierającym $I$.
Równość $I^S=I$ nie musi zachodzić.
\end{twierdzenie}
\begin{proof}
Z Lematu~\ref{aa:lem:local-ideal-membership} wynika,
że $a/1\in S^{-1}I$ dokładnie wtedy, gdy $sa\in I$
dla pewnego $s\in S$. To jest dokładnie definicja
przeciwobrazu oraz zbioru $I^S$.
Przeciwobraz ideału przez homomorfizm jest ideałem:
jeśli $a,b\in I^S$, to $sa,tb\in I$ dla pewnych
$s,t\in S$, zatem $st(a-b)=t(sa)-s(tb)\in I$;
ponadto $s(ra)=r(sa)\in I$ dla każdego $r\in R$.
Wreszcie $I\subseteq I^S$, gdyż $1\in S$.
\end{proof}

\begin{przyklad}[Nasycenie może powiększyć ideał]
\label{aa:prz:localization-saturation-z}
Niech $R=\mathbb Z$, $S=\{2^n:n\ge0\}$
oraz $I=(6)$. Mamy $I^S=(3)$.
Jeśli $a\in I^S$, to dla pewnego $n$ liczba $6$
dzieli $2^na$, co wymusza $3\mid a$.
Odwrotnie, gdy $3\mid a$, liczba $2a$ jest
podzielna przez $6$, więc $a\in I^S$.
W lokalizacji $\mathbb Z[1/2]$ ideały generowane
przez $6$ i przez $3$ są jednak identyczne,
ponieważ $6=2\cdot3$ i $2$ stało się jednostką.
\end{przyklad}

\begin{twierdzenie}[Pierwiastnik ideału a lokalizacja]
\label{aa:tw:localization-radical}
Niech $R$ będzie przemiennym pierścieniem z jedynką,
$S\subseteq R$ zbiorem multiplikatywnym
niezawierającym zera, a $I\subseteq R$ ideałem.
W pierścieniu $S^{-1}R$ zachodzi równość ideałów
\[
 \sqrt{S^{-1}I}=S^{-1}(\sqrt I).
\]
\end{twierdzenie}
\begin{proof}
Weźmy $a/s\in S^{-1}(\sqrt I)$.
Z definicji rozszerzenia ideału istnieją
$b\in\sqrt I$ i $t\in S$ takie, że $a/s=b/t$.
Ponieważ $b\in\sqrt I$, istnieje liczba naturalna
$n\ge1$ z $b^n\in I$. Stąd
$(a/s)^n=b^n/t^n\in S^{-1}I$, czyli
$a/s\in\sqrt{S^{-1}I}$. To daje pierwsze zawieranie.

Dla zawierania odwrotnego niech
$a/s\in\sqrt{S^{-1}I}$. Dla pewnego $n\ge1$
mamy $a^n/s^n\in S^{-1}I$.
Na mocy Lematu~\ref{aa:lem:local-ideal-membership}
istnieje $u\in S$ takie, że $ua^n\in I$.
Ponieważ $I$ jest ideałem, także
$(ua)^n=u^{n-1}(ua^n)\in I$.
Z definicji pierwiastnika wynika $ua\in\sqrt I$.
Równość ułamków $a/s=(ua)/(us)$ jest poprawna,
gdyż oba iloczyny krzyżowe są równe $aus$.
Ponieważ $ua\in\sqrt I$ i $us\in S$,
otrzymujemy $a/s\in S^{-1}(\sqrt I)$.
Równość ideałów została udowodniona.
\end{proof}

\begin{przyklad}[Pierwiastniki po odwróceniu dwójki]
\label{aa:prz:localization-radical-example}
Niech $R=\mathbb Z$, $I=(12)$ oraz
$S=\{2^n:n\in\mathbb N_0\}$.
W pierścieniu $\mathbb Z$ zachodzi
$\sqrt{(12)}=(6)$: liczba całkowita ma pewną
potęgę podzielną przez $12=2^2\cdot3$ wtedy
i tylko wtedy, gdy sama jest podzielna przez
$2$ i $3$. Po odwróceniu liczby $2$ mamy
$S^{-1}(12)=(3)$ i $S^{-1}(6)=(3)$,
gdzie ideały zapisano już w $\mathbb Z[1/2]$.
Istotnie, $12=3\cdot 2^2$ i $6=3\cdot2$,
a potęgi $2$ są jednostkami lokalizacji.
Twierdzenie daje więc
$\sqrt{(3)}=(3)$ w $\mathbb Z[1/2]$,
co zgadza się także z faktem, że
$\mathbb Z[1/2]/(3)\cong\mathbb F_3$ jest ciałem.
\end{przyklad}

\section{Ideały pierwsze po lokalizacji i pierścienie lokalne}
\begin{twierdzenie}[Odpowiedniość ideałów pierwszych przy lokalizacji]
\label{aa:tw:prime-local-correspondence}
Niech $R$ będzie przemiennym pierścieniem z jedynką,
a $S\subseteq R$ zbiorem multiplikatywnym niezawierającym zera.
Niech $\iota:R\to S^{-1}R$ oznacza homomorfizm
kanoniczny $\iota(a)=a/1_R$. Przyporządkowania
$P\mapsto S^{-1}P$ oraz $Q\mapsto\iota^{-1}(Q)$
są wzajemnie odwrotnymi
bijekcjami między ideałami pierwszymi $P$ pierścienia $R$
spełniającymi $P\cap S=\varnothing$ a ideałami
pierwszymi $Q$ lokalizacji $S^{-1}R$.
\end{twierdzenie}
\begin{proof}
Niech $P$ będzie ideałem pierwszym pierścienia $R$
rozłącznym z $S$. Rozszerzenie $S^{-1}P$ jest właściwe
na mocy Lematu~\ref{aa:lem:local-ideal-membership}.
Jeśli iloczyn $(a/s)(b/t)=ab/(st)$ należy do tego
ideału, istnieje $u\in S$ z $uab\in P$.
Ponieważ $u\notin P$ i $P$ jest pierwszy, otrzymujemy
$ab\in P$, a następnie $a\in P$ lub $b\in P$.
Zatem przynajmniej jeden z ułamków $a/s$, $b/t$
należy do $S^{-1}P$, co dowodzi jego pierwszości.
Ponadto $a/1_R\in S^{-1}P$ dokładnie wtedy,
gdy $ua\in P$ dla pewnego $u\in S$.
Pierwszość $P$ i $u\notin P$ dają wtedy $a\in P$.
W ten sposób przeciwobraz $S^{-1}P$ jest równy $P$.

Odwrotnie, niech $Q$ będzie ideałem pierwszym $S^{-1}R$
i połóżmy $P=\iota^{-1}(Q)$. Przeciwobraz ideału
pierwszego przez homomorfizm jest pierwszy:
$ab\in P$ oznacza $\iota(a)\iota(b)\in Q$, więc
$\iota(a)\in Q$ lub $\iota(b)\in Q$.
Ponieważ elementy $\iota(s)$ dla $s\in S$ są jednostkami,
nie mogą należeć do właściwego ideału $Q$, zatem
$P\cap S=\varnothing$.
Wreszcie dla każdego $a/s\in S^{-1}R$ zachodzi
\[
 \frac as\in Q
 \quad\Longleftrightarrow\quad
 \frac a1=\frac s1\frac as\in Q
 \quad\Longleftrightarrow\quad a\in P.
\]
W pierwszej równoważności wykorzystaliśmy odwracalność
$s/1$: z przynależności $a/1$ do ideału wynika
przynależność $(s/1)^{-1}(a/1)=a/s$.
Stąd $Q=S^{-1}P$, a obie konstrukcje są odwrotne.
\end{proof}

\begin{definicja}[Pierścień lokalny]
\label{aa:def:local-ring}
Niech $R$ będzie przemiennym, niezerowym pierścieniem
z jedynką. Pierścień $R$ nazywamy \emph{lokalnym},
jeśli ma dokładnie jeden ideał maksymalny.
Dla ideału pierwszego $P\subsetneq R$ definiujemy
\emph{lokalizację w ideale pierwszym}
$R_P=(R\setminus P)^{-1}R$.
Zbiór $R\setminus P$ jest multiplikatywny:
$1_R\notin P$, a jeśli $a,b\notin P$, to
$ab\notin P$ na mocy pierwszości $P$.
\end{definicja}

\begin{twierdzenie}[Lokalizacja w ideale pierwszym jest pierścieniem lokalnym]
\label{aa:tw:localization-at-prime}
Niech $R$ będzie przemiennym pierścieniem z jedynką,
a $P\subsetneq R$ ideałem pierwszym. Wtedy
$R_P=(R\setminus P)^{-1}R$ jest pierścieniem lokalnym,
a jego jedynym ideałem maksymalnym jest
$PR_P=(R\setminus P)^{-1}P$.
Element $a/s\in R_P$ jest jednostką wtedy i tylko wtedy,
gdy $a\notin P$.
\end{twierdzenie}
\begin{proof}
Dla $s\notin P$ element $s/1$ jest jednostką.
Jeśli $a\notin P$, także $a/1$ jest jednostką,
a więc $a/s=(a/1)(s/1)^{-1}$ jest jednostką.
Jeżeli $a\in P$, to $a/s\in PR_P$, które jest
ideałem właściwym na mocy rozłączności
$P\cap(R\setminus P)=\varnothing$.
Żadna jednostka nie leży w ideale właściwym,
więc taki ułamek nie jest jednostką.

Udowodniliśmy, że zbiór elementów nieodwracalnych
pierścienia $R_P$ jest dokładnie $PR_P$.
Jeśli $J\subsetneq R_P$ jest dowolnym ideałem właściwym,
nie może zawierać jednostki, więc wszystkie jego
elementy należą do $PR_P$. Stąd $J\subseteq PR_P$.
Wynika z tego, że $PR_P$ jest maksymalny i że każdy
inny ideał maksymalny musi być równy $PR_P$.
\end{proof}

\begin{twierdzenie}[Ciało reszt lokalizacji]
\label{aa:tw:localization-residue-field}
Niech $R$ będzie przemiennym pierścieniem z jedynką,
a $P\subsetneq R$ jego ideałem pierwszym. Niech
$R_P=(R\setminus P)^{-1}R$ i $PR_P=(R\setminus P)^{-1}P$.
Wówczas
\[
 R_P/PR_P\ \cong\ \operatorname{Frac}(R/P).
\]
Po prawej stronie występuje ciało ułamków dziedziny
całkowitości $R/P$.
\end{twierdzenie}
\begin{proof}
Ponieważ $P$ jest pierwszy, $R/P$ jest dziedziną,
a więc ma ciało ułamków opisane w
Lemacie~\ref{aa:lem:fracfield-well-defined}.
Dla $a\in R$ i $s\in R\setminus P$ określmy
\[
 \Phi(a/s)=(a+P)/(s+P)\in\operatorname{Frac}(R/P).
\]
Mianownik jest niezerowy. Jeśli $a/s=b/t$,
to dla pewnego $u\notin P$ mamy $u(at-bs)=0$.
Po przejściu do $R/P$ i skróceniu niezerowej klasy
$u+P$ otrzymujemy $(a+P)(t+P)=(b+P)(s+P)$;
zatem $\Phi$ jest dobrze określone. Ze wzorów na
ułamki wynika, że jest homomorfizmem pierścieni.
Każdy element $\operatorname{Frac}(R/P)$ ma postać
$(a+P)/(s+P)$ z $s\notin P$, więc $\Phi$ jest
surjektywny. Ułamek $a/s$ należy do jądra dokładnie
wtedy, gdy $a+P=0$, czyli $a\in P$; to z kolei
jest równoważne $a/s\in PR_P$ na mocy
Twierdzenia~\ref{aa:tw:localization-at-prime}.
Pierwsze twierdzenie o izomorfizmie kończy dowód.
\end{proof}

Poniższy diagram komutatywny przedstawia jednocześnie
lokalizację, przejście do ilorazu i zanurzenie w ciało
ułamków. Strzałka $\Phi$ ma jądro $PR_P$, a dolna
strzałka jest zanurzeniem dziedziny $R/P$ w jej ciało
ułamków.
\begin{center}
\begin{tikzpicture}[>=Stealth]
\node (R) at (0,1.8) {$R$};
\node (Loc) at (5.0,1.8) {$R_P$};
\node (Quot) at (0,0) {$R/P$};
\node (Frac) at (5.0,0) {$\operatorname{Frac}(R/P)$};
\draw[->] (R) -- node[above] {$\iota$} (Loc);
\draw[->] (R) -- node[left] {$\pi$} (Quot);
\draw[->] (Loc) -- node[right] {$\Phi$} (Frac);
\draw[->] (Quot) -- node[below] {$\jmath$} (Frac);
\end{tikzpicture}
\end{center}
W obu drogach z $R$ do $\operatorname{Frac}(R/P)$
element $a$ przechodzi w $(a+P)/(1+P)$.

\begin{przyklad}[Pierścień liczb wymiernych o mianownikach niepodzielnych przez $5$]
\label{aa:prz:z-local-five}
Dla $R=\mathbb Z$ i $P=(5)$ mamy
$\mathbb Z_{(5)}=\{a/b:\ a,b\in\mathbb Z,\ 5\nmid b\}\subset\mathbb Q$.
Na przykład $7/12$ i $1/3$ należą do tego pierścienia,
a $1/5$ nie należy: gdyby $1/5=a/b$ z $5\nmid b$,
z równości $b=5a$ wynikałaby sprzeczność.
Element $a/b$ jest jednostką dokładnie wtedy,
gdy $5\nmid a$, bo wówczas jego odwrotność $b/a$
również ma mianownik niepodzielny przez $5$.
Jedynym ideałem maksymalnym jest
$5\mathbb Z_{(5)}$, a ciało reszt jest izomorficzne
z $\mathbb F_5$ poprzez $a/b\mapsto[a][b]^{-1}$.
\end{przyklad}

\begin{przyklad}[Lokalizacja wielomianów w ideale $(x)$]
\label{aa:prz:local-kxy}
Weźmy $R=K[x,y]$ nad ciałem $K$ oraz
$P=(x)$. Jest to ideał pierwszy, bo
$K[x,y]/(x)\cong K[y]$ jest dziedziną całkowitości.
Pierścień $R_P$ zawiera odwrotności wszystkich
wielomianów spoza $(x)$, czyli tych, których
wartość po podstawieniu $x=0$ jest niezerowym
wielomianem zmiennej $y$.
W szczególności $y$ staje się jednostką w $R_P$,
ponieważ $y\notin(x)$, ale $x$ pozostaje nieodwracalny.
Jedynym ideałem maksymalnym $R_P$ jest $(x)R_P$.
Nie należy mylić tej konstrukcji z $K[x,x^{-1}]$,
w którym odwracamy potęgi $x$. W tym pierścieniu
$x$ jest jednostką, ale pierścień nie jest lokalny.
Jeśli ciało $K$ ma co najmniej trzy elementy,
wybierzmy $c\in K\setminus\{0,1\}$. Ideały $(x-1)$
i $(x-c)$ są różne i maksymalne, ponieważ podstawienie
odpowiednio $x=1$ i $x=c$ daje surjekcje
$K[x,x^{-1}]\to K$ o tych właśnie jądrach.
Jeśli $K=\mathbb F_2$, można zamiast drugiego użyć
ideału $(x^2+x+1)$. Wielomian $p=x^2+x+1$ nie ma
pierwiastka w $\mathbb F_2$, bo dla $x=0$ i $x=1$
przyjmuje wartość $1$; jako wielomian stopnia dwa jest
zatem nierozkładalny. W ciele $\mathbb F_2[x]/(p)$
zachodzi $[x]([x]+[1])=[1]$, gdyż
$x(x+1)=p-1=p+1$ w charakterystyce $2$.
Stąd klasa $x$ jest odwracalna, a więc
$\mathbb F_2[x,x^{-1}]/(p)\cong\mathbb F_2[x]/(p)$
jest ciałem. Ideał $(p)$ w pierścieniu Laurenta
jest więc maksymalny i różny od $(x-1)$, ponieważ
$p(1)=1\ne0$.
Istotna różnica jest następująca:
\emph{odwrócenie jednego elementu nie oznacza lokalizacji
w jednym ideale pierwszym}.
\end{przyklad}

\begin{uwaga}[Sprostowanie dotyczące pierścieni z dzielnikami zera]
\label{aa:uw:localization-zerodivisors}
W pierścieniu z dzielnikami zera lokalizacja może
zidentyfikować dwa różne elementy. Na przykład
w $R=\mathbb Z/6\mathbb Z$ dla
$S=\{[1],[2],[4]\}$ mamy $[3]\cdot[2]=[0]$,
więc klasa $[3]/[1]$ w $S^{-1}R$ jest zerowa.
Lokalizacja jest tu izomorficzna z $\mathbb F_3$:
redukcja modulo $3$ wysyła każdy element $S$ na
jednostkę $\mathbb F_3$, więc z własności uniwersalnej
otrzymujemy surjektywny homomorfizm
$S^{-1}R\to\mathbb F_3$. Jeżeli $[a]/[2^n]$ ma obraz
zerowy, to $3\mid a$; wtedy $[a]$ w $R$ jest klasą
$[0]$ lub $[3]$, a obie stają się zerem w lokalizacji.
Homomorfizm jest więc także różnowartościowy.
Nie wolno zatem w ogólnym przypadku zastępować relacji
z Definicji~\ref{aa:def:localization} przez zwykłą
równość iloczynów krzyżowych bez dodatkowego czynnika.
\end{uwaga}

\section{Skończone generowanie ideałów i przygotowanie do teorii Noethera}
\begin{definicja}[Pierścień noetherowski]
\label{aa:def:noetherian-ring-intro}
Przemienny pierścień z jedynką $R$ jest
\emph{noetherowski}, jeśli każdy jego ideał jest
skończenie generowany, to znaczy ma postać
$(a_1,\ldots,a_r)$ dla pewnego $r\in\mathbb N_0$
i pewnych $a_1,\ldots,a_r\in R$; dla $r=0$
przyjmujemy ideał zerowy $(0)$.
\end{definicja}

\begin{twierdzenie}[Warunek łańcucha i skończone generowanie ideałów]
\label{aa:tw:noetherian-acc}
Dla przemiennego pierścienia $R$ z jedynką
następujące warunki są równoważne:
\begin{enumerate}[label=\textup{(\roman*)}]
\item każdy ideał $R$ jest skończenie generowany;
\item każdy wstępujący ciąg ideałów
$I_1\subseteq I_2\subseteq\cdots$ od pewnego miejsca jest stały.
\end{enumerate}
\end{twierdzenie}
\begin{proof}
Załóżmy (i) i rozważmy wstępujący ciąg ideałów.
Jego suma mnogościowa $I=\bigcup_{n\ge1}I_n$ jest
ideałem: jeśli $a\in I_j$ i $b\in I_k$, to
$a-b\in I_{\max(j,k)}$, a dla dowolnego $r\in R$
mamy $ra\in I_j$. Istnieją generatory
$a_1,\ldots,a_m$ ideału $I$, gdzie $m\ge0$.
Jeśli $m=0$, to $I=(0)$, a ponieważ każdy
$I_n\subseteq I$, cały ciąg składa się z ideałów
zerowych i jest stały. Załóżmy więc $m\ge1$.
Każdy generator $a_i$ należy do pewnego $I_{n_i}$;
dla $N=\max\{n_1,\ldots,n_m\}$ wszystkie generatory
leżą w $I_N$. Stąd $I=(a_1,\ldots,a_m)\subseteq I_N$.
Odwrotne zawieranie jest oczywiste, zatem $I_N=I$
i dla wszystkich $n\ge N$ mamy $I_n=I_N$.

Załóżmy (ii). Przypuśćmy, że istnieje ideał $I$,
który nie jest skończenie generowany. Wybierzmy
$a_1\in I$ i ideał $I_1=(a_1)$ różny od $I$.
Jest to możliwe, bo $I$ nie jest generowany przez
jeden element. Następnie wybierzmy
$a_2\in I\setminus I_1$ i połóżmy
$I_2=(a_1,a_2)$. Ponieważ $I$ nie jest skończenie
generowany, $I_2\ne I$ i można wybrać
$a_3\in I\setminus I_2$. Powtarzając konstrukcję,
otrzymamy ściśle rosnący ciąg
$I_1\subsetneq I_2\subsetneq I_3\subsetneq\cdots$,
sprzeczny z (ii). Zatem wszystkie ideały są
skończenie generowane.
\end{proof}

\begin{przyklad}[Pierścienie noetherowskie i nie-noetherowskie]
\label{aa:prz:noetherian-example}
Niech $K$ będzie ciałem. Każda dziedzina ideałów
głównych jest noetherowska, ponieważ każdy jej ideał
ma jeden generator. W szczególności $\mathbb Z$
i $K[x]$ są noetherowskie.

Pierścień $K[x_1,x_2,\ldots]$ wielomianów w przeliczalnie
wielu zmiennych nie jest noetherowski. Istotnie, ciąg
ideałów $I_n=(x_1,\ldots,x_n)$ jest ściśle rosnący.
Gdyby $x_{n+1}\in I_n$, istniałyby wielomiany $f_i$
takie, że $x_{n+1}=\sum_{i=1}^n x_if_i$.
Po podstawieniu $x_1=\cdots=x_n=0$ prawa strona
wynosiłaby zero, a lewa $x_{n+1}$, co jest sprzecznością
w pierścieniu wielomianów pozostałych zmiennych.
\end{przyklad}

\begin{uwaga}[Dalsza kolejność pojęć]
\label{aa:uw:ring-further}
W kolejnym etapie rozbudujemy teorię pierścieni
noetherowskich i artinowskich, twierdzenie Hilberta
o bazie, rozkłady pierwotne ideałów i ich zastosowania.
Wcześniejsze wyniki o pierwiastnikach, lokalizacji,
ideałach pierwszych i maksymalnych stanowią do tego
bezpośrednie przygotowanie; nie wymagają pełnej
teorii modułów.
\end{uwaga}

\setlength{\parskip}{4pt}
